% =====================================================================
%  Short paper (maximo 5 paginas) -- formato de conferencia IEEE (IEEEtran)
%  PPI Grupo 4 -- UPeU Juliaca -- protocolo de deteccion de anomalias PeMS-D4
%  Version en espanol del 27 oct 2026 (fin de la semana 8 de 12).
%
%  LEYENDA DE ESTADO usada en el texto:
%    "obtenido"  -> resultados producidos y verificados con pruebas
%                   automaticas (assert) (Anexo C del PPI, semanas 1--8).
%    "pendiente" -> analisis definidos a priori pero aun no ejecutados.
%  Las lineas marcadas  %% TODO  requieren informacion de los autores.
% =====================================================================
\documentclass[conference]{IEEEtran}
\usepackage[utf8]{inputenc}
\usepackage[T1]{fontenc}
\usepackage[spanish,es-nodecimaldot]{babel}
\usepackage{cite}
\usepackage{amsmath,amssymb,amsfonts}
\usepackage{graphicx}
\usepackage{textcomp}
\usepackage{xcolor}
\usepackage{booktabs}
\usepackage{pifont}

\newcommand{\cm}{\ding{51}}          % reportado / cumplido
\newcommand{\xm}{\ding{55}}          % no realizado
\newcommand{\nr}{n.i.}               % no informado
\newcommand{\pend}{\textit{pendiente}} % analisis aun no ejecutado
\newcommand{\pl}{\textit{plan}}            % planificado, aun no ejecutado
\usepackage{url}

\begin{document}

\title{Aprendizaje profundo frente a una l\'inea base no param\'etrica para la
detecci\'on de anomal\'ias en el flujo vehicular de autopistas: protocolo
uniforme y dise\~no de sensibilidad al etiquetado en PeMS-D4}

%% TODO: confirmar orden de autores, correos u ORCID, y si el asesor figura
%% como coautor segun las reglas de la conferencia.
\author{\IEEEauthorblockN{Wilbert Alex Mayta Arotaype}
\IEEEauthorblockA{\textit{Escuela Profesional de Ingenier\'ia de Sistemas} \\
\textit{Universidad Peruana Uni\'on}\\
Juliaca, Per\'u \\
correo o ORCID}
\and
\IEEEauthorblockN{Juan Enders Prieto Quispe}
\IEEEauthorblockA{\textit{Escuela Profesional de Ingenier\'ia de Sistemas} \\
\textit{Universidad Peruana Uni\'on}\\
Juliaca, Per\'u \\
correo o ORCID}
\and
\IEEEauthorblockN{Wrangler Choquechambi Luque}
\IEEEauthorblockA{\textit{Escuela Profesional de Ingenier\'ia de Sistemas} \\
\textit{Universidad Peruana Uni\'on}\\
Juliaca, Per\'u \\
correo o ORCID}
\and
\IEEEauthorblockN{Esteban Tocto Cano}
\IEEEauthorblockA{\textit{Escuela Profesional de Ingenier\'ia de Sistemas} \\
\textit{Universidad Peruana Uni\'on}\\
Juliaca, Per\'u \\
correo o ORCID}
}

\maketitle

% NOTA (regla de la plantilla): sin simbolos, caracteres especiales, notas
% al pie ni matematicas en el titulo ni en el resumen.
\begin{abstract}
Los detectores de anomal\'ias de tr\'afico basados en aprendizaje profundo
suelen reportarse con conjuntos de datos, particiones, umbrales y etiquetas
heterog\'eneos, mientras que estudios recientes muestran que los detectores
estad\'isticos simples suelen ser competitivos una vez eliminados los
defectos de evaluaci\'on. Presentamos un protocolo uniforme, cuyas reglas
principales se fijaron antes de ver cualquier resultado de redes neuronales,
que compara seis familias de aprendizaje profundo (LSTM, GRU, CNN-1D,
autoencoder, Transformer e h\'ibrido CNN-LSTM-GRU) y un modelo publicado
(SI-DMSeq) frente a una l\'inea base no param\'etrica de estad\'isticos m\'oviles
sobre series de flujo de autopistas de PeMS-D4. El protocolo usa validaci\'on
temporal progresiva, un procedimiento com\'un y reducido de ajuste para las
seis familias entrenadas (el modelo publicado queda en su configuraci\'on por
defecto), tres semillas para los modelos finales, PR-AUC sin ajuste por
puntos, pruebas pareadas con correcci\'on de Holm, medici\'on de latencia y dos
esquemas de etiquetado construidos por separado: un proxy estad\'istico causal
y la inyecci\'on controlada de anomal\'ias puntuales, contextuales y
colectivas, para evaluar si el ranking de detectores depende de c\'omo se
definen las anomal\'ias. Los resultados preliminares, que pasaron las
verificaciones autom\'aticas de consistencia especificadas, abarcan la
depuraci\'on de datos (291 de 307 sensores conservados), ambos esquemas de
etiquetado y la l\'inea base. El esquema proxy produce una prevalencia de 0.10
por ciento y solo 29.2 por ciento de los bloques de evaluaci\'on contiene
ambas clases; su conjunto positivo casi no se solapa con el inyectado
(\'indice de Jaccard de a lo sumo 0.0012); y la l\'inea base alcanza un PR-AUC
medio por bloque de 0.345 en prueba bajo inyecci\'on, unas 11 veces el azar,
pero un PR-AUC agregado de solo 0.003 bajo el proxy, unas tres veces el azar.
Los tiempos de entrenamiento medidos obligaron a una b\'usqueda de
hiperpar\'ametros reducida y declarada de antemano. Los resultados de
aprendizaje profundo y las pruebas estad\'isticas est\'an pendientes.
\end{abstract}

\begin{IEEEkeywords}
detecci\'on de anomal\'ias, aprendizaje profundo, protocolo de evaluaci\'on,
PeMS, series temporales, flujo vehicular, inyecci\'on sint\'etica de anomal\'ias
\end{IEEEkeywords}

% ---------------------------------------------------------------------
\section{Introducci\'on}

\textbf{Contexto y afectados.} Los centros de gesti\'on de tr\'afico, como el
Departamento de Transporte de California (Caltrans), que patrocina PeMS
\cite{chen2003pems}, dependen de sensores de autopista para detectar
perturbaciones, y los conductores sufren las demoras. La Administraci\'on
Federal de Carreteras de EE.\ UU.\ atribuye cerca de la mitad de la
congesti\'on a perturbaciones temporales, y solo los incidentes explican 25\%
de la congesti\'on \cite{fhwa_nonrecurring}. En el \'area de la Bah\'ia de San
Francisco, INRIX estim\'o que la congesti\'on cost\'o a cada conductor 79 horas y
USD~2\,250 en 2017, el a\~no previo al periodo estudiado aqu\'i
\cite{cbs2018inrix}. Detectar a tiempo observaciones at\'ipicas tiene por tanto
valor operativo. Si los detectores se adoptan por resultados que no se pueden
comparar, un operador puede pagar modelos profundos costosos sin una ventaja
demostrada sobre una regla simple, y por eso hace falta un protocolo uniforme
de evaluaci\'on.

Los detectores de lazo en autopistas producen series de flujo densas,
estacionales y ruidosas, en las que identificar a tiempo observaciones
at\'ipicas apoya el monitoreo del tr\'afico. Se han propuesto detectores de
aprendizaje profundo (DL) basados en recurrencia, convoluci\'on,
reconstrucci\'on o atenci\'on
\cite{zhang2022automatic,mendez2022using,ijaradar2025datadriven,nouri2025sidmseq},
pero sus resultados son dif\'iciles de comparar: los estudios difieren en
datos, particiones, etiquetas, umbrales y m\'etricas, y los grandes conjuntos
p\'ublicos de autopistas como PeMS se construyeron para pron\'ostico y no
traen etiquetas de anomal\'ia \cite{coursey2024ftaed}. Por ello las etiquetas
se derivan de bit\'acoras de incidentes, reglas o simulaci\'on
\cite{sun2022datacentric,yang2025sumo}, y una anomal\'ia estad\'istica no es
necesariamente un incidente real.

Dos l\'ineas de evidencia motivan una evaluaci\'on m\'as estricta. El ajuste por
puntos (PA) infla tanto los puntajes que puntajes aleatorios alcanzan F1 de
estado del arte \cite{kim2022rigorous}, y los benchmarks populares son en
gran parte resolubles con reglas de una l\'inea basadas en estad\'isticos m\'oviles
\cite{wu2022flawed}. Cuando la evaluaci\'on es controlada, los detectores
cl\'asicos igualan o superan a los profundos en benchmarks gen\'ericos
\cite{schmidl2022anomaly,audibert2022deep,rewicki2023worth,liu2024elephant}
y en volumen de tr\'afico etiquetado por expertos
\cite{taheri2026applicability}; adem\'as, la calidad de las etiquetas cambia
las conclusiones en la detecci\'on de incidentes viales
\cite{sun2022datacentric}, y los rankings var\'ian con el tipo de anomal\'ia y
la m\'etrica \cite{liu2024elephant,sorbo2024navigating}.

\textbf{Objeto de investigaci\'on y problema.} El objeto de estudio es la
detecci\'on de anomal\'ias en el flujo vehicular de autopistas con aprendizaje
profundo, analizada en PeMS-D4 con el bloque detector $\times$ esquema de
etiquetado $\times$ sensor-pliegue como unidad de an\'alisis. El problema es si
los detectores DL ofrecen una ventaja reproducible sobre una l\'inea base de
estad\'isticos m\'oviles, y si el ranking resultante se mantiene estable ante
distintas definiciones operativas de anomal\'ia y bajo una restricci\'on de
costo de inferencia.

\textbf{Vac\'io de conocimiento.} Ya existen comparaciones uniformes para series
gen\'ericas \cite{schmidl2022anomaly,liu2024elephant} y para volumen de tr\'afico
urbano \cite{taheri2026applicability}. Lo que no se sabe es si la posici\'on
relativa de los detectores DL frente a una l\'inea base simple en flujo de
autopistas sobrevive a un cambio en la forma de definir las anomal\'ias. En la
literatura revisada, ning\'un estudio examina esto sobre una misma partici\'on
temporal de un conjunto de autopistas y a la vez compara familias DL
gen\'ericas, una l\'inea base no param\'etrica y un modelo de tr\'afico publicado,
trata la latencia de inferencia como criterio de primer orden y cuantifica la
variabilidad entre semillas con control de comparaciones m\'ultiples
(Tabla~\ref{tab:practices}). El vac\'io es por tanto de evidencia sobre
robustez, no de detectores.

\textbf{Objetivos y preguntas.} \emph{Objetivo general:} determinar si los
detectores DL ofrecen una ventaja reproducible sobre una l\'inea base de
estad\'isticos m\'oviles (G0) en PeMS-D4 y si el ranking de detectores es
estable entre esquemas de etiquetado. \emph{Objetivos espec\'ificos:}
(O1)~estimar, bajo el Esquema B, la diferencia de PR-AUC entre cada detector
DL y G0; (O2)~establecer cu\'ales de los detectores que superan a G0 cumplen un
criterio de latencia predefinido; (O3)~cuantificar la concordancia entre los
rankings de detectores bajo los Esquemas A y B. \emph{Preguntas de
investigaci\'on:} RQ1~(O1, H1), \textquestiondown cu\'anto mejora, si es que lo
hace, alg\'un detector DL la precisi\'on promedio de G0 bajo el Esquema B;
RQ2~(O2, H2), \textquestiondown cumplen el criterio de latencia los detectores
que mejoran a G0; RQ3~(O3, H3), \textquestiondown es estable el ranking de
detectores entre esquemas de etiquetado? An\'alisis con entradas controladas y
por tipo de anomal\'ia prueban la robustez de estas comparaciones.

\textbf{Aportes declarados.} (1)~\emph{Emp\'irico:} evidencia sobre si la
posici\'on de los detectores DL frente a una l\'inea base m\'ovil en flujo de
autopistas es robusta a la definici\'on de anomal\'ia (resultados DL
pendientes); (2)~\emph{metodol\'ogico:} un protocolo reproducible para PeMS-D4
que operacionaliza pr\'acticas destiladas de los estudios semilla, con dos
esquemas de etiquetado construidos por separado y un criterio a~priori de
estabilidad del ranking; (3)~\emph{pr\'actico:} la latencia como criterio
expl\'icito, con un frente de Pareto PR-AUC$\times$latencia. Los resultados
preliminares aqu\'i reportados conciernen solo a los datos, las etiquetas y la
l\'inea base.

% ---------------------------------------------------------------------
\section{Trabajos relacionados y pr\'acticas destiladas}

\textbf{Detecci\'on de anomal\'ias de tr\'afico.} Los detectores reportados
incluyen pron\'ostico con atenci\'on sobre grafos \cite{zhang2022automatic},
autoencoders CNN y BiLSTM \cite{mendez2022using}, un autoencoder
GCN-LSTM-atenci\'on evaluado con anomal\'ias artificiales
\cite{ijaradar2025datadriven} y SI-DMSeq, un codificador--decodificador LSTM
multitarea que marca perturbaciones con una carta de control
\cite{nouri2025sidmseq}. El etiquetado es el principal obst\'aculo: la
supervisi\'on d\'ebil elev\'o el F1 de la clase incidente de un LSTM de 0.69 a
0.88 en el Distrito~4 de PeMS \cite{sun2022datacentric}, y FT-AED combina
registros de choques con etiquetas de expertos \cite{coursey2024ftaed}. En
datos de volumen etiquetados por expertos, un perfil horario
media$\pm c\cdot$desv. alcanz\'o F1 0.566 frente a 0.120 de un LSTM
\cite{taheri2026applicability}.

\textbf{Defectos de evaluaci\'on.} Con PA, puntajes uniformemente aleatorios
alcanzan F1 0.969 en SWaT frente a 0.216 sin \'el \cite{kim2022rigorous}; 316
de 367 series de Yahoo son separables con una l\'inea de media/desviaci\'on
m\'oviles \cite{wu2022flawed}. El AUC-ROC est\'a dominado por los verdaderos
negativos bajo desbalance \cite{sorbo2024navigating,liu2024elephant}, pero el
PR-AUC puntual no es robusto a peque\~nos desfases temporales
\cite{liu2024elephant}.

\textbf{Etiquetas sint\'eticas.} Las anomal\'ias se dividen en puntuales
(globales, contextuales) y por patr\'on \cite{lai2021revisiting}. Las anomal\'ias
inyectadas ayudan a entrenar y seleccionar detectores, pero difieren de las
reales y solo transfieren cuando coinciden los tipos
\cite{carmona2022ncad,goswami2023unsupervised}.

La Tabla~\ref{tab:practices} resume las caracter\'isticas de protocolo
reportadas por los estudios semilla y motiva cada decisi\'on de dise\~no de la
Secci\'on~III.

\begin{table}[t]
\centering
\caption{Caracter\'isticas de protocolo reportadas por los estudios semilla.
Tr: datos de tr\'afico; NB: l\'inea base no profunda; CS: partici\'on cronol\'ogica
declarada; NT: sin PA ni umbral ajustado en prueba para la m\'etrica
principal; RR: corridas repetidas; ST: prueba de significancia; LT:
tiempo/latencia. \cm\ s\'i, \xm\ no, p parcial, \nr\ no informado, -- no
aplica, plan: planificado, a\'un no ejecutado.}
\label{tab:practices}
\setlength{\tabcolsep}{3.2pt}
\renewcommand{\arraystretch}{0.9}
\footnotesize
\begin{tabular}{@{}lccccccc@{}}
\toprule
Estudio & Tr & NB & CS & NT & RR & ST & LT \\
\midrule
Taheri'26 \cite{taheri2026applicability} & \cm & \cm & \nr & \xm & \nr & \xm & \xm \\
Taheri'24 \cite{taheri2024tpfds}         & \cm & \cm & \nr & \xm & \cm & \xm & \xm \\
FT-AED'24 \cite{coursey2024ftaed}        & \cm & p   & \nr & \cm & \xm & \xm & \xm \\
Sun'22 \cite{sun2022datacentric}         & \cm & \xm & \cm & \cm & \cm & \xm & \xm \\
Zhang'22 \cite{zhang2022automatic}       & \cm & \cm & \nr & \cm & \xm & \xm & \xm \\
Ijaradar'25 \cite{ijaradar2025datadriven}& \cm & \cm & \nr & \nr & \nr & \xm & \xm \\
Nouri'25 \cite{nouri2025sidmseq}         & \cm & \xm & \cm & --  & \xm & \xm & \cm \\
Kim'22 \cite{kim2022rigorous}            & \xm & \cm & --  & \xm & p   & \xm & \xm \\
Audibert'22 \cite{audibert2022deep}      & \xm & \cm & --  & \xm & \nr & \cm & \xm \\
Schmidl'22 \cite{schmidl2022anomaly}     & \xm & \cm & --  & \cm & \nr & \xm & \cm \\
Liu'24 \cite{liu2024elephant}            & \xm & \cm & --  & \cm & \xm & \cm & \nr \\
\midrule
\textbf{Este protocolo}                  & \cm & \cm & \cm & \cm & \pl & \pl & \pl \\
\bottomrule
\end{tabular}
\end{table}

% ---------------------------------------------------------------------
\section{Protocolo}

\textbf{Dise\~no.} Experimento computacional comparativo con estructura
factorial (detector $\times$ esquema de etiquetado) y medidas repetidas sobre
bloques sensor$\times$pliegue, con validaci\'on temporal progresiva y un bloque
de prueba reservado.

\textbf{Datos.} PeMS-D4 \cite{guo2019astgcn} contiene flujo, ocupaci\'on y
velocidad de 307 sensores de autopista del \'area de la Bah\'ia de San
Francisco, del 1~ene al 28~feb 2018, con resoluci\'on de 5~min (16\,992
pasos por sensor). Se excluyeron \emph{a priori} los sensores con m\'as de 20\%
de ceros o una racha congelada de 24~h, quedando 291.
%% TODO: confirmar que este control de calidad examino todo el periodo antes
%% de la particion; si es asi, indicarlo aqui (un revisor pidio la declaracion).
Cada muestra es una ventana de 12 pasos ($t-11..t$, una hora) con $F=8$
entradas (flujo, ocupaci\'on, velocidad, codificaciones c\'iclicas de hora y
d\'ia de la semana, y un indicador de hora pico); la etiqueta es la del paso
final $t$, de modo que todos los m\'etodos punt\'uan el mismo instante. Se entrena
un modelo global por arquitectura con las ventanas de todos los sensores.

\textbf{Partici\'on.} El \'ultimo 20\% de cada serie es un bloque de prueba
reservado que se usa una sola vez. El primer 80\% forma tres pliegues
progresivos de validaci\'on temporal; la normalizaci\'on z-score, los
hiperpar\'ametros y los umbrales se ajustan solo con las partes de
entrenamiento y validaci\'on de cada pliegue, nunca con datos de prueba
(Tabla~\ref{tab:practices}, CS/NT).

\textbf{Esquemas de etiquetado.} El \emph{Esquema A} (proxy) etiqueta
$x_{s,t}$ como an\'omalo cuando $|z_{s,t}|>3$, con media y desviaci\'on
est\'andar calculadas causalmente sobre los 288 pasos previos de flujo crudo; la prevalencia
anticipada era de 0.3--2\%.
La l\'inea base G0 pertenece a la misma familia de reglas (media y desviaci\'on
m\'oviles, aunque con una ventana mucho m\'as corta), por lo que A se usa solo
para la prueba de estabilidad del ranking (H3) y nunca para H1. El
\emph{Esquema B} inyecta anomal\'ias puntuales, contextuales y colectivas (3,
6 o 12 pasos) en series reales, perturbando flujo, ocupaci\'on y velocidad,
con 3\% de prevalencia objetivo por partici\'on y por bloque sensor$\times$
pliegue, semilla de inyecci\'on fija, restricciones f\'isicas (flujo y
velocidad $\geq 0$; ocupaci\'on en $[0,1]$) y sin solapamiento entre
particiones. Las anomal\'ias contextuales reemplazan un valor por uno menor al
percentil 5 o mayor al 95 de un perfil hora$\times$d\'ia de la semana por
sensor, estimado con datos de entrenamiento, una familia de reglas distinta de
G0 \cite{lai2021revisiting}.

\textbf{Condiciones.} G0: detector m\'ovil no param\'etrico solo sobre flujo;
punt\'ua cada paso con el z-score absoluto sobre los $w$ pasos previos (ventana
$w\in\{6,11,22,44\}$ pasos y multiplicador de umbral $k\in\{2.0,2.5,3.0\}$,
una grilla de 12 puntos). G1:
LSTM; G2: GRU; G3: CNN-1D; G4: autoencoder (error de reconstrucci\'on como
puntaje); G5: encoder Transformer; G6: h\'ibrido CNN-LSTM-GRU; G7: SI-DMSeq
\cite{nouri2025sidmseq} con su configuraci\'on publicada, usando el error de
reconstrucci\'on local de la serie en el paso $t$ como puntaje puntual. G1--G3,
G5 y G6 optimizan una p\'erdida supervisada de clasificaci\'on (salida sigmoide,
entrop\'ia cruzada binaria). G4 optimiza una p\'erdida de reconstrucci\'on sobre
ventanas cuya etiqueta del paso final es normal bajo el esquema
correspondiente, por lo que esas ventanas pueden contener pasos an\'omalos
anteriores. G0 y G7 no optimizan ninguna p\'erdida supervisada de
clasificaci\'on, aunque se usa selecci\'on basada en etiquetas donde se
especifica (ventana y umbral de G0 en validaci\'on). Arquitectura y grado de
uso de etiquetas est\'an por tanto parcialmente confundidos, y la comparaci\'on
es entre sistemas detectores y no entre arquitecturas aisladas. G0--G6 se
implementan en TensorFlow/Keras. Cada una de G1--G6 eval\'ua 12 configuraciones
por esquema, la primera es su especificaci\'on inicial y las otras 11 se
sortean de los rangos fijos de la Tabla~\ref{tab:ranges} con semilla
determinista. G7 se mantiene en PyTorch y
sin ajuste, una asimetr\'ia declarada a~priori.

\textbf{B\'usqueda reducida.} La b\'usqueda nominal (6 arquitecturas $\times$ 2
esquemas $\times$ 3 pliegues $\times$ 12 configuraciones $\times$ 3 semillas)
son 1\,296 entrenamientos; con los 4.3--5.4 min por \'epoca medidos en el
pliegue 3 (CPU, LSTM), 30 \'epocas cada uno superan 70 d\'ias. Antes de
cualquier b\'usqueda fijamos una reducci\'on para todas las arquitecturas: una
semilla, como m\'aximo 10 \'epocas con parada temprana (paciencia 3) sobre
PR-AUC, una de cada cuatro ventanas de entrenamiento y selecci\'on sobre el
bloque de validaci\'on \emph{completo}. La ganadora se reentrena con todo el
bloque de 80\%, tres semillas, todas las ventanas y sin parada temprana,
usando la mediana de la mejor \'epoca. El costo se cronometr\'o en un AMD Ryzen 7 7445HS
(6 n\'ucleos, 12 hilos) con 15.2\,GB de RAM, en CPU (TensorFlow 2.16.2); una
corrida en GPU (RTX 4050, bajo WSL2) fue m\'as lenta para este modelo y tama\~no
de lote (unos 12 frente a 5.4 min por \'epoca). Cada corrida registra su
hardware y los tiempos de entrenamiento se comparan entre arquitecturas solo si
se midieron en la misma m\'aquina.
%% TODO: agregar la maquina usada para latencia si es distinta de la anterior.
%% TODO: confirmar que las Decisiones 15--17 estan aprobadas por el equipo y
%% que config_final.json esta sellado antes del envio (sellar_config_final.py).
%% TODO: los rangos de busqueda de abajo son propuesta (Decision 17); confirmar.

\begin{table}[t]
\centering
\caption{Espacio de b\'usqueda de G1--G6. El primer valor de cada lista es la
especificaci\'on inicial (configuraci\'on 0); las otras 11 configuraciones se
sortean de estas listas.}
\label{tab:ranges}
\footnotesize
\setlength{\tabcolsep}{3pt}
\renewcommand{\arraystretch}{0.9}
\begin{tabular}{@{}p{0.2\columnwidth}p{0.72\columnwidth}@{}}
\toprule
Modelo & Valores \\
\midrule
G1 LSTM, G2 GRU & unidades $\{[64,32],[32,16],[128,64]\}$; dropout $\{0.1,0.2,0.3\}$; tasa de aprendizaje $\{5\!\times\!10^{-4},10^{-3},2\!\times\!10^{-3}\}$ \\
G3 CNN-1D & filtros $\{[32,64,128],[16,32,64],[64,128,256]\}$; kernel $\{3,5\}$; decaimiento de pesos $\{10^{-4},0,10^{-3}\}$; tasa como G1 \\
G4 autoencoder & capas $\{[32,16,8],[64,32,16],[32,16,4]\}$; L2 $\{0,10^{-5},10^{-4}\}$; tasa como G1 \\
G5 Transformer & bloques $\{2,1,3\}$; cabezas 4; $d_{model}$ $\{64,32\}$; feed-forward $\{128,64\}$; dropout $\{0.1,0.2\}$; tasa $\{5\!\times\!10^{-4},3\!\times\!10^{-4},10^{-3}\}$; decaimiento $\{10^{-4},10^{-3}\}$ \\
G6 CNN-LSTM-GRU & filtros $\{64,32\}$; kernel $\{3,5\}$; unidades LSTM $\{64,32\}$; unidades GRU $\{32,16\}$; dropout $\{0.2,0.1,0.3\}$; tasa como G1 \\
\bottomrule
\end{tabular}
\end{table}

\textbf{Selecci\'on y m\'etricas.} Se elige la configuraci\'on que maximiza el
PR-AUC medio de validaci\'on sobre los tres pliegues; el umbral de F1 y el
umbral que alcanza recall 0.90 se eligen en el \'ultimo bloque de validaci\'on
(los modelos finales se reentrenan con ese bloque, lo que puede sesgar F1 y
FPR pero no el PR-AUC). La m\'etrica primaria es el PR-AUC (precisi\'on
promedio, AP), cuyo nivel de azar es la prevalencia. El AP medio por bloque es
la media simple del AP por sensor sobre los bloques evaluables; el AP agregado
se calcula sobre puntajes y etiquetas concatenados de todos los pasos v\'alidos
de una partici\'on y se reporta con su prevalencia. El PR-AUC de G0 depende
solo de $w$; $k$ fija el umbral binario usado para F1. F1, AUC-ROC y la tasa
de falsos positivos (FPR) en el umbral de F1 y en el umbral de recall 0.90 son
secundarios, y el recall realmente alcanzado en prueba se reporta junto con
esa FPR; el PA est\'a prohibido \cite{kim2022rigorous,sorbo2024navigating}.
Como las anomal\'ias colectivas dominan los instantes inyectados, el AP se
reporta tambi\'en por tipo de anomal\'ia, tomando como positivos los instantes
de un tipo y como negativos los instantes normales (los otros tipos se
excluyen). La latencia se mide con \texttt{perf\_counter} tras 100 llamadas de
calentamiento sobre $\geq$1000 llamadas, con lote 1 y una ventana por sensor,
reportando mediana, p95 e IC 95\%; el criterio operativo es $\leq$10\,ms por
ventana, complementado con el frente de Pareto PR-AUC$\times$latencia.

\textbf{Controles contra un atajo de firma.} Las redes ven ocho canales, G0
solo el flujo, y el Esquema B perturba varios canales, por lo que una red
podr\'ia explotar la firma de la inyecci\'on en lugar de detectar mejor las
anomal\'ias de flujo. El AP por tipo y una ablaci\'on reentrenada con ocupaci\'on
y velocidad en cero lo abordan; la ablaci\'on conserva los canales de hora,
d\'ia de la semana y hora pico, por lo que la informaci\'on sigue sin igualar a
la de G0. El entrenamiento con pesos de clase prueba la sensibilidad al
desbalance, no el atajo. Estos an\'alisis se fijaron tras una corrida
exploratoria m\'inima que no se reporta como resultado.

\textbf{Hip\'otesis y estad\'istica.} H1: al menos un detector DL supera a G0
en PR-AUC medio por bloque en prueba bajo el Esquema B por al menos
$\delta=0.05$ de PR-AUC absoluto (cerca de 14\% del 0.345 de G0), un margen que
descarta ganancias triviales y se fija antes de abrir el bloque de prueba (H0:
ninguno lo hace). H2, condicionada a H1: al menos
un detector que cumple H1 satisface el criterio de latencia (mediana de una
ventana, lote 1, solo inferencia del modelo, CPU com\'un; tambi\'en se reporta
p95); H0: ninguno. H3: el ranking de detectores no es estable entre esquemas
($\tau$ de Kendall $<0.7$); H0: $\tau\geq0.7$. Cada alternativa enuncia el efecto a demostrar (una ventaja,
cumplir la latencia, un cambio de ranking), por lo que H3 se formula como
inestabilidad. Los pliegues de validaci\'on
sirven para selecci\'on y comparaciones exploratorias (prueba de Friedman sobre
los $291\times3$ bloques, Nemenyi \cite{demsar2006statistical}); H1 se eval\'ua
primordialmente sobre los 291 bloques de prueba reservados, despu\'es de
congelar configuraciones, semillas, reglas de agregaci\'on y comparaciones, con
pruebas pareadas frente a G0 y correcci\'on de Holm ($\alpha=0.05$). Sensores y
pliegues progresivos no son independientes y Holm no corrige la dependencia,
por lo que la significancia es indicativa y se reportan tama\~nos de efecto e
intervalos. El AP por bloque se calcula por semilla y se promedia sobre las
tres semillas antes de la prueba; tambi\'en se reportan resultados por semilla.
Los bloques sin alguna de las clases se excluyen y se reportan. Para H3 las
condiciones se ordenan por PR-AUC medio por bloque en cada esquema (bloques
v\'alidos en ambos, remuestreados para un IC 95\% bootstrap de $\tau$): el
ranking es estable si todo el IC est\'a en 0.7 o por encima, H3 se apoya si
est\'a por debajo, y el resultado es inconcluso en otro caso; la $\rho$ de
Spearman es un an\'alisis de sensibilidad. Como H3 cambia a la vez etiquetas,
prevalencia, ejemplos de entrenamiento, configuraciones seleccionadas y datos
perturbados, eval\'ua la sensibilidad del procedimiento completo a la
definici\'on de anomal\'ia, no solo de las etiquetas de evaluaci\'on. Bajo el
Esquema A la ventaja sobre G0 se reporta solo de forma descriptiva, porque A
comparte la familia de reglas de G0.
%% TODO (decisiones del equipo antes de abrir el bloque de prueba para las
%% redes): (a) registrar como decision nueva que H1 es primaria en prueba,
%% pues el PPI aun indica pliegues de validacion; (b) nombrar la prueba
%% pareada (p. ej., Wilcoxon de rangos con signo); (c) confirmar media sobre
%% semillas (el PPI dice media o mediana); (d) confirmar mediana vs p95 en H2;
%% (e) confirmar el margen
%% delta = 0.05 propuesto aqui (cerca de 14% del 0.345 de G0) o reemplazarlo,
%% antes de abrir el bloque de prueba; (f) indicar bases de datos, cadena de
%% busqueda, fechas y criterios de inclusion de la revision de literatura que
%% respalda "en la literatura revisada".

% ---------------------------------------------------------------------
\section{Resultados preliminares}

La Tabla~\ref{tab:results} reporta la evidencia obtenida hasta ahora
(semanas 1--8) y lo que sigue pendiente. El pipeline, ambos esquemas de
etiquetado y G0 pasaron 50 de 50 verificaciones autom\'aticas, incluidas
pruebas de perturbaci\'on que alteran valores futuros; el cargador de ventanas
pas\'o 73 de 73 y las definiciones de modelos 54 de 54 verificaciones frente a
la especificaci\'on inicial.

\begin{table}[t]
\centering
\caption{Evidencia obtenida a la fecha frente a an\'alisis pendientes. El
nivel de azar es igual a la prevalencia.}
\label{tab:results}
\footnotesize
\setlength{\tabcolsep}{3pt}
\renewcommand{\arraystretch}{0.9}
\begin{tabular}{@{}p{0.47\columnwidth}p{0.45\columnwidth}@{}}
\toprule
\'Item & Resultado \\
\midrule
\multicolumn{2}{@{}l}{\textit{Obtenido}}\\
Sensores conservados tras control de calidad & 291 de 307 (16 excluidos, 5.2\%) \\
Prevalencia Esquema A & 0.10\% (4\,873 / 4\,860\,864); 0.086--0.125\% por partici\'on \\
Prevalencia Esquema B & 3.02--3.06\% por bloque \\
Composici\'on Esquema B (instantes) & 74.7\% colectiva, 14.6\% puntual, 10.7\% contextual \\
Solapamiento A vs.\ B (Jaccard) & 0.0003--0.0012 \\
Bloques con ambas clases & A: 340/1\,164 (29.2\%); B: 100\% \\
Config.\ elegida de G0 (A y B) & ventana 22, multiplicador 3.0 (multiplicador en el borde de la grilla; no afecta el PR-AUC) \\
G0 prueba, Esquema B & PR-AUC medio por bloque 0.345 ($\approx$11.3$\times$ azar); agregado 0.335 (11.0$\times$); F1 0.323; AUC-ROC 0.792; FPR 54.3\% en el umbral que alcanza recall 0.90 en validaci\'on \\
G0 prueba, Esquema B, por tipo & puntual 0.905; contextual 0.399; colectiva 0.131 \\
G0 prueba, Esquema A & PR-AUC medio por bloque 0.084 (84 de 291 bloques evaluables); agregado 0.003 (3.0$\times$); F1 0.021; FPR 65.6\% en el umbral de recall 0.90 derivado de validaci\'on \\
\midrule
\multicolumn{2}{@{}l}{\textit{Pendiente}}\\
B\'usqueda y entrenamiento final G1--G6 & \pend \\
Friedman/Holm, Nemenyi (H1) & \pend \\
Latencia, frente de Pareto (H2); prueba de humo G7 & \pend \\
$\tau$ de Kendall, $\rho$ de Spearman (H3) & \pend \\
Ablaci\'on de canales, pesos de clase & \pend \\
\bottomrule
\end{tabular}
\end{table}

%% TODO: agregar el recall alcanzado en prueba con el umbral de recall 0.90
%% derivado de validacion (lo imprime "recall prueba" semana6_g0_ampliado.py).
Cuatro observaciones orientan el trabajo restante. Primero, las etiquetas
proxy e inyectadas son casi disjuntas (Jaccard $\leq0.0012$), por lo que H3 no
se cumple de forma trivial. Segundo, el Esquema A tiene una prevalencia menor
al 0.3--2\% anticipado y solo 29.2\% de bloques con ambas clases, lo que limita
la intersecci\'on de bloques sobre la que se calcula H3; adem\'as sus dos
agregaciones discrepan para G0 (0.084 medio por bloque sobre 84 bloques
evaluables, 0.003 agregado, unas tres veces la prevalencia agregada), por lo
que cada PR-AUC se reporta con su agregaci\'on. Tercero, G0 est\'a muy por
encima del azar bajo inyecci\'on, en l\'inea con la competitividad de los
detectores simples \cite{taheri2026applicability,schmidl2022anomaly}, pero
depende mucho del tipo: separa casi perfectamente las anomal\'ias puntuales y
mal las colectivas, que son tres cuartas partes de los instantes inyectados,
por lo que cualquier ganancia de DL debe leerse por tipo. Cuarto, el costo de
entrenamiento medido (4.3--5.4 min por \'epoca en CPU) oblig\'o a la b\'usqueda
reducida de la Secci\'on~III, solo parcialmente mitigada por el reentrenamiento
con varias semillas.

% ---------------------------------------------------------------------
\section{Discusi\'on y amenazas a la validez}

La evidencia previa de que las l\'ineas base estad\'isticas pueden rivalizar con
DL \cite{audibert2022deep,taheri2026applicability} motiva H1, pero no detectar
una ventaja de DL se reportar\'a con tama\~nos de efecto e incertidumbre, sin
interpretar la no significancia como evidencia de equivalencia. La latencia
publicada de SI-DMSeq (unos 55\,ms por secuencia para 46 zonas y ventanas de
24 pasos \cite{nouri2025sidmseq}) corresponde a otra carga, por lo que todos
los m\'etodos se medir\'an bajo una misma configuraci\'on; el criterio de 10\,ms
(unos 3\,s para una ventana por sensor, muy por debajo del intervalo de
muestreo de 5~min) es un criterio t\'ecnico de este estudio.

\emph{Validez externa:} un solo conjunto de datos, una red de autopistas y unos
once d\'ias de prueba limitan la generalizaci\'on; H1 concierne a anomal\'ias
inyectadas, no a incidentes verificados. \emph{Validez de constructo:} ning\'un esquema representa incidentes
verificados; las anomal\'ias inyectadas difieren de las reales
\cite{carmona2022ncad} y los eventos reales sin etiquetar act\'uan como
pseudo-negativos \cite{goswami2023unsupervised}. \emph{Validez interna:} el
Esquema A comparte familia de reglas con G0, de ah\'i su rol restringido; todo
ajuste es local al pliegue. G0 usa solo flujo y una ventana de 22 pasos frente
a 12 pasos y ocho canales de las redes, y la b\'usqueda reducida de una semilla
puede favorecer arquitecturas que convergen temprano. Los controles de atajo
se fijaron tras una peque\~na corrida exploratoria. \emph{Validez de
conclusi\'on:} los pliegues progresivos no son independientes y los sensores
comparten din\'amica, por lo que el an\'alisis de Friedman en validaci\'on es
exploratorio y se planea un an\'alisis solo del pliegue 3 como chequeo de
robustez. \emph{Validez de la m\'etrica:} el PR-AUC puntual da m\'as peso a las
anomal\'ias largas (colectivas) y es sensible al desfase
\cite{sorbo2024navigating,liu2024elephant}; se reportar\'an resultados por
tipo. \emph{Validez de implementaci\'on:} G7 corre en otro framework, un
confusor declarado para la latencia; el hardware de entrenamiento y de latencia
se documentar\'a por separado.

% ---------------------------------------------------------------------
\section{Conclusi\'on}

Especificamos y ejecutamos parcialmente un protocolo uniforme y reproducible
para comparar detectores DL con una l\'inea base no param\'etrica sobre flujo de
autopistas de PeMS-D4 bajo dos esquemas de etiquetado construidos por
separado. Los resultados preliminares muestran que los esquemas producen
conjuntos positivos en gran medida no solapados, que la l\'inea base est\'a muy
por encima del azar bajo inyecci\'on pero depende del tipo de anomal\'ia, y que
el presupuesto nominal de ajuste era inviable, lo que llev\'o a una b\'usqueda
reducida declarada. Siguen la bater\'ia DL, las pruebas de H1--H3, los
controles de atajo y la medici\'on de latencia; el c\'odigo (licencia MIT), las
configuraciones y las semillas se publicar\'an con los resultados.
%% TODO: agregar la URL del repositorio cuando sea publico.

\bibliographystyle{IEEEtran}
\bibliography{references}

\end{document}
